\chapter{Funzioni}\label{cap:funzioni}

Le funzioni sono particolari corrispondenze: quelle che associano a \textbf{ogni} elemento del primo insieme
\textbf{uno e un solo} elemento del secondo. In questo capitolo le definiamo, le componiamo e studiamo le tre
proprietà fondamentali (iniettività, suriettività, biettività), caratterizzandole in più modi.

% ──────────────────────────────────────────────────────────────
\section{Definizione di funzione}\label{sec:def-funzione}

\dfn[funzione]{Funzione}{
    Siano $a$ e $b$ insiemi. Una corrispondenza $f = (a, b, g)$ si dice \textbf{applicazione} o \textbf{funzione} se
    \[ \forall x \in a\ \exists!\, y \in b\ (x f y). \]
    Visto che $y$ è unico, lo denoto $f(x)$. Inoltre:
    \begin{itemize}
        \item $f(x)$ si dice \textbf{immagine di $x$ mediante $f$};
        \item $a$ e $b$ si dicono rispettivamente \textbf{dominio} e \textbf{codominio} di $f$;
        \item $\mathrm{Im} f := \{y \in b \mid \exists x \in a\,(f(x) = y)\} = \{f(x) \mid x \in a\}$ si dice \textbf{immagine di $f$};
        \item scriviamo brevemente $f \colon a \to b$.
    \end{itemize}
}

Nella scrittura $\{f(x) \mid x \in a\}$ non si indica l'insieme di appartenenza perché è implicitamente il codominio di $f$.
È a volte utile dare una \textbf{descrizione esplicita} della funzione, che evidenzi per ogni $x$ qual è la sua immagine:
\[ f \colon x \in a \mapsto f(x) \in b. \]

% ──────────────────────────────────────────────────────────────
\section{Composizione e funzioni notevoli}\label{sec:composizione}

\begin{concetto}
\begin{Proposizione}{Funzione composta}{composta}
    \begin{prerequisiti}
        \dip{Definizione}{def:funzione}
        \dip{Definizione}{def:prodotto-corrispondenze}
    \end{prerequisiti}
    Siano $f \colon a \to b$ e $g \colon b \to c$ funzioni. Il prodotto $gf$ (come corrispondenze) è una funzione da $a$ a $c$,
    detta \textbf{funzione composta} e indicata con $g \circ f = gf$, con descrizione esplicita
    \[ g \circ f \colon x \in a \mapsto g(f(x)) \in c. \]
\end{Proposizione}
\begin{dimostrazione}
    Sia $x \in a$. \textbf{Esistenza}: posto $y = f(x) \in b$ e $z = g(y) \in c$, si ha $x f y$ e $y g z$, quindi
    $x \,(gf)\, z$ per definizione di prodotto. \textbf{Unicità}: se $x \,(gf)\, z'$, esiste $w$ con $x f w$ e $w g z'$;
    per l'unicità in $f$ si ha $w = f(x)$, e per l'unicità in $g$ si ha $z' = g(w) = g(f(x))$.
    Quindi ogni $x \in a$ ha un'unica immagine, che è $g(f(x))$.
\end{dimostrazione}
\end{concetto}

\dfn[funzioni-notevoli]{Funzioni notevoli}{
    \begin{itemize}
        \item Se $c \in b$, la funzione $f \colon x \in a \mapsto c \in b$ la dico \textbf{funzione costante} (di valore $c$).
        \item Se $a \subseteq b$, la funzione $f \colon x \in a \mapsto x \in b$ la dico \textbf{immersione} di $a$ in $b$.
        \item Se $s \subseteq a$, la funzione $f_{|s} \colon x \in s \mapsto f(x) \in b$ si chiama \textbf{restrizione} di $f$ a $s$,
              e $f$ si dice \textbf{prolungamento} di $f_{|s}$ ad $a$.
        \item La funzione $x \in a \mapsto x \in a$ la dico \textbf{identità} (o funzione identica) su $a$ e la indico con $\mathrm{id}_a$.
    \end{itemize}
}

% ──────────────────────────────────────────────────────────────
\section{Iniettività, suriettività, biettività}\label{sec:iniettive}

\dfn[iniettiva]{Funzione iniettiva}{
    $f \colon a \to b$ si dice \textbf{iniettiva} se
    \[ \forall x_1, x_2 \in a\,\big(x_1 \neq x_2 \rightarrow f(x_1) \neq f(x_2)\big). \]
}

\begin{concetto}
\begin{Proposizione}{Altre forme dell'iniettività}{forme-iniettivita}
    \begin{prerequisiti}
        \dip{Definizione}{def:iniettiva}
        \dip{Proposizione}{prop:contrapposizione}
        \dip{Teorema}{teo:negazione-quantificatori}
        \dip{Proposizione}{prop:negimp}
        \dip{Proposizione}{prop:doppianeg}
    \end{prerequisiti}
    Sia $f \colon a \to b$.
    \begin{enumerate}
        \item $f$ è iniettiva se e solo se $\forall x_1, x_2 \in a\,\big(f(x_1) = f(x_2) \rightarrow x_1 = x_2\big)$.
        \item $f$ \textbf{non} è iniettiva se e solo se $\exists x_1, x_2 \in a\,\big(x_1 \neq x_2 \wedge f(x_1) = f(x_2)\big)$.
    \end{enumerate}
\end{Proposizione}
\begin{dimostrazione}
    \textbf{1.} Per contrapposizione $(x_1 \neq x_2 \rightarrow f(x_1) \neq f(x_2)) \leftrightarrow (\neg(f(x_1) \neq f(x_2)) \rightarrow \neg(x_1 \neq x_2))$,
    e per la doppia negazione questo è $f(x_1) = f(x_2) \rightarrow x_1 = x_2$.

    \textbf{2.} Negando la definizione: $\neg\forall x_1, x_2\,(\dots)$ diventa $\exists x_1, x_2\,\neg(x_1 \neq x_2 \rightarrow f(x_1) \neq f(x_2))$
    (negazione dei quantificatori), cioè $\exists x_1, x_2\,(x_1 \neq x_2 \wedge f(x_1) = f(x_2))$ (negazione dell'implicazione e doppia negazione).
\end{dimostrazione}
\end{concetto}

La forma 1 è quella più comoda nelle dimostrazioni: \textbf{per dimostrare che $f$ è iniettiva si prendono $x_1, x_2$ con
$f(x_1) = f(x_2)$ e si mostra che $x_1 = x_2$}.

\dfn[suriettiva]{Funzione suriettiva}{
    $f \colon a \to b$ si dice \textbf{suriettiva} se
    \[ \forall y \in b\ \exists x \in a\,\big(f(x) = y\big). \]
    Se $\mathrm{Im} f \subseteq t \subseteq b$, la funzione $x \in a \mapsto f(x) \in t$ si chiama \textbf{ridotta} di $f$ a $t$.
}

\begin{concetto}
\begin{Proposizione}{La ridotta all'immagine è suriettiva}{ridotta}
    \begin{prerequisiti}
        \dip{Definizione}{def:suriettiva}
        \dip{Definizione}{def:funzione}
    \end{prerequisiti}
    Sia $f \colon a \to b$. La ridotta di $f$ a $\mathrm{Im} f$, cioè $x \in a \mapsto f(x) \in \mathrm{Im} f$, è suriettiva.
\end{Proposizione}
\begin{dimostrazione}
    Sia $y \in \mathrm{Im} f$. Per definizione di immagine esiste $x \in a$ con $f(x) = y$.
\end{dimostrazione}
\end{concetto}

\dfn[biettiva]{Funzione biettiva}{
    Se $f$ è suriettiva e iniettiva, $f$ si dice \textbf{biettiva} (o \textbf{biunivoca}).
}

\begin{concetto}
\begin{Teorema}{Composta di funzioni iniettive}{comp-iniettive}
    \begin{prerequisiti}
        \dip{Definizione}{def:iniettiva}
        \dip{Proposizione}{prop:forme-iniettivita}
        \dip{Proposizione}{prop:composta}
    \end{prerequisiti}
    Siano $f \colon a \to b$ e $g \colon b \to c$ iniettive. Allora $g \circ f$ è iniettiva.
\end{Teorema}
\begin{dimostrazione}
    Siano $x_1, x_2 \in a$ tali che $g \circ f(x_1) = g \circ f(x_2)$, cioè $g(f(x_1)) = g(f(x_2))$.
    Ma $g$ è iniettiva, allora $f(x_1) = f(x_2)$; ma $f$ è iniettiva, allora $x_1 = x_2$.
\end{dimostrazione}
\end{concetto}

\begin{concetto}
\begin{Teorema}{Composta di funzioni suriettive}{comp-suriettive}
    \begin{prerequisiti}
        \dip{Definizione}{def:suriettiva}
        \dip{Proposizione}{prop:composta}
    \end{prerequisiti}
    Siano $f \colon a \to b$ e $g \colon b \to c$ suriettive. Allora $g \circ f$ è suriettiva.
\end{Teorema}
\begin{dimostrazione}
    Sia $y \in c$. Poiché $g$ è suriettiva, trovo $z \in b$ con $g(z) = y$; ma $f$ è suriettiva, allora trovo
    $x \in a$ con $f(x) = z$. Quindi $g(f(x)) = g(z) = y$.
\end{dimostrazione}
\end{concetto}

\begin{concetto}
\begin{Corollario}{Composta di funzioni biettive}{comp-biettive}
    \begin{prerequisiti}
        \dip{Teorema}{teo:comp-iniettive}
        \dip{Teorema}{teo:comp-suriettive}
        \dip{Definizione}{def:biettiva}
    \end{prerequisiti}
    La composta di funzioni biettive è biettiva.
\end{Corollario}
\begin{dimostrazione}
    Segue immediatamente dai due teoremi precedenti.
\end{dimostrazione}
\end{concetto}

\begin{concetto}
\begin{Teorema}{Funzioni composte e caratteristiche}{composte-caratteristiche}
    \begin{prerequisiti}
        \dip{Proposizione}{prop:composta}
        \dip{Proposizione}{prop:forme-iniettivita}
        \dip{Definizione}{def:suriettiva}
    \end{prerequisiti}
    Siano $f \colon a \to b$ e $g \colon b \to c$ due funzioni.
    \begin{enumerate}
        \item Se $g \circ f$ è iniettiva, allora $f$ è iniettiva.
        \item Se $g \circ f$ è suriettiva, allora $g$ è suriettiva.
    \end{enumerate}
    Cioè, se $g \circ f$ è biettiva, allora $f$ è iniettiva e $g$ è suriettiva.
\end{Teorema}
\begin{dimostrazione}
    \textbf{1.} Prendo $x_1, x_2 \in a$ con $f(x_1) = f(x_2)$; quindi $g \circ f(x_1) = g(f(x_1)) = g(f(x_2)) = g \circ f(x_2)$,
    ma $g \circ f$ è iniettiva, quindi $x_1 = x_2$. Dunque $f$ è iniettiva.

    \textbf{2.} Sia $y \in c$. Visto che $g \circ f$ è suriettiva e il suo codominio è $c$, trovo $x \in a$ con
    $g \circ f(x) = y$, cioè $g(f(x)) = y$: l'elemento $f(x) \in b$ ha immagine $y$ mediante $g$. Questo vale per ogni $y$,
    quindi $g$ è suriettiva.
\end{dimostrazione}
\end{concetto}

\begin{esempio}{Tre esempi importanti}
    \begin{itemize}
        \item Sia $a$ un insieme. La funzione $x \in a \mapsto \{x\} \in \mcP(a)$ è \textbf{iniettiva} (se $\{x\} = \{y\}$, per
              Estensionalità $x = y$) ma \textbf{non suriettiva}: $\varnothing \in \mcP(a)$ non è un singleton, quindi non è immagine di nessun $x$.
        \item La funzione $x \in \mcP(a) \mapsto a \setminus x \in \mcP(a)$ è una importantissima funzione \textbf{biettiva}.
              Infatti per la doppia negazione insiemistica (Proposizione~\ref{prop:doppia-neg-ins}) $a \setminus (a \setminus x) = x$
              per ogni $x \subseteq a$: se $a \setminus x = a \setminus y$, applicando $a \setminus$ a entrambi si ottiene $x = y$
              (iniettiva); e ogni $y \subseteq a$ è l'immagine di $a \setminus y$ (suriettiva).
        \item Siano $a = \{1\}$, $b = \{1, 2\}$, $f$ l'immersione di $a$ in $b$ e $g$ la funzione costante da $b$ in $a$.
              Allora $g \circ f = \mathrm{id}_a$, ma $f$ non è suriettiva e $g$ non è iniettiva: nel Teorema~\ref{teo:composte-caratteristiche}
              non si può dire di più.
    \end{itemize}
    \begin{center}
    \begin{tikzpicture}[>=Stealth, every node/.style={font=\small}]
        \node (a1) at (0,0) {$1$};
        \node (b1) at (3,0.5) {$1$}; \node (b2) at (3,-0.5) {$2$};
        \draw[rounded corners] (-0.5,-1) rectangle (0.5,1); \node at (0,1.3) {$a = \{1\}$};
        \draw[rounded corners] (2.5,-1) rectangle (3.5,1); \node at (3,1.3) {$b = \{1,2\}$};
        \draw[->,primary] (a1) to[bend left=15] node[above,sloped] {$f$} (b1);
        \draw[->,cherryred] (b1) to[bend left=15] node[below,sloped] {$g$} (a1);
        \draw[->,cherryred] (b2) to[bend left=10] (a1);
    \end{tikzpicture}
    \end{center}
\end{esempio}

% ──────────────────────────────────────────────────────────────
\section{Immagine e antiimmagine}\label{sec:immagine}

\dfn[immagine-antiimmagine]{Applicazioni immagine e antiimmagine}{
    Sia $f \colon a \to b$. Definisco l'applicazione \textbf{immagine}
    \[ \overrightarrow{f} \colon x \in \mcP(a) \mapsto \overrightarrow{f}(x) = \{f(z) \mid z \in x\} \in \mcP(b) \]
    e l'applicazione \textbf{antiimmagine}
    \[ \overleftarrow{f} \colon y \in \mcP(b) \mapsto \overleftarrow{f}(y) = \{x \in a \mid f(x) \in y\} \in \mcP(a). \]
}

\begin{concetto}
\begin{Proposizione}{Proprietà di immagine e antiimmagine}{prop-immagine}
    \begin{prerequisiti}
        \dip{Definizione}{def:immagine-antiimmagine}
        \dip{Definizione}{def:funzione}
        \dip{Definizione}{def:suriettiva}
        \dip{Lemma}{lem:vuoto}
        \dip{Proposizione}{prop:doppia-inclusione}
    \end{prerequisiti}
    Sia $f \colon a \to b$ e sia $x \in a$. Allora:
    \begin{enumerate}
        \item $\overrightarrow{f}(\varnothing) = \varnothing = \overleftarrow{f}(\varnothing)$;
        \item $\overrightarrow{f}(a) = \mathrm{Im} f$;
        \item $f$ è suriettiva $\leftrightarrow$ $\mathrm{Im} f = b$ $\leftrightarrow$ $\overrightarrow{f}(a) = b$;
        \item $\overrightarrow{f}(\{x\}) = \{f(x)\}$;
        \item $\overleftarrow{f}(b) = a = \overleftarrow{f}(\mathrm{Im} f)$.
    \end{enumerate}
\end{Proposizione}
\begin{dimostrazione}
    \textbf{1.} Non esiste $z \in \varnothing$, quindi $\{f(z) \mid z \in \varnothing\}$ non ha elementi; allo stesso modo nessun
    $x$ soddisfa $f(x) \in \varnothing$. Per il Lemma~\ref{lem:vuoto} entrambi sono $\varnothing$.

    \textbf{2.} $\overrightarrow{f}(a) = \{f(z) \mid z \in a\}$, che è la definizione di $\mathrm{Im} f$.

    \textbf{3.} Sempre $\mathrm{Im} f \subseteq b$; la suriettività dice esattamente che ogni $y \in b$ sta in $\mathrm{Im} f$, cioè
    $b \subseteq \mathrm{Im} f$. Per la doppia inclusione, $f$ suriettiva $\leftrightarrow \mathrm{Im} f = b$; l'ultima equivalenza segue dal punto 2.

    \textbf{4.} $\overrightarrow{f}(\{x\}) = \{f(z) \mid z \in \{x\}\} = \{f(z) \mid z = x\} = \{f(x)\}$.

    \textbf{5.} Per ogni $x \in a$ si ha $f(x) \in \mathrm{Im} f \subseteq b$, quindi ogni $x \in a$ sta sia in $\overleftarrow{f}(b)$ sia in
    $\overleftarrow{f}(\mathrm{Im} f)$; viceversa entrambi sono per definizione sottoinsiemi di $a$. Per doppia inclusione sono uguali ad $a$.
\end{dimostrazione}
\end{concetto}

% ──────────────────────────────────────────────────────────────
\section{Caratterizzazioni con l'antiimmagine}\label{sec:caratt-antiimmagine}

\dfn[map-in-su-bi]{Insiemi di funzioni}{
    Siano $a$, $b$ due insiemi. Dico
    \begin{itemize}
        \item $\mathrm{Map}(a, b)$ l'insieme delle funzioni da $a$ a $b$;
        \item $\mathrm{In}(a, b)$ l'insieme delle funzioni iniettive da $a$ a $b$;
        \item $\mathrm{Su}(a, b)$ l'insieme delle funzioni suriettive da $a$ a $b$;
        \item $\mathrm{Bi}(a, b)$ l'insieme delle funzioni biettive da $a$ a $b$.
    \end{itemize}
}

\begin{concetto}
\begin{Teorema}{Caratterizzazione delle funzioni iniettive}{caratt-iniettive}
    \begin{prerequisiti}
        \dip{Definizione}{def:map-in-su-bi}
        \dip{Definizione}{def:immagine-antiimmagine}
        \dip{Proposizione}{prop:forme-iniettivita}
        \dip{Proposizione}{prop:negdoppiaimp}
        \dip{Teorema}{teo:negazione-quantificatori}
        \dip{Definizione}{def:esiste-unico}
    \end{prerequisiti}
    $f \colon a \to b$ è iniettiva se e solo se, per ogni $y \in b$, $\overleftarrow{f}(\{y\})$ ha al più un elemento, cioè:
    \[ \forall a, b\ \forall f\,\Big(f \in \mathrm{In}(a, b) \leftrightarrow \forall y \in b\,\big(\overleftarrow{f}(\{y\}) = \varnothing \vee \exists!\, x \in a\,(f(x) = y)\big)\Big). \]
\end{Teorema}
\begin{dimostrazione}
    Lo dimostro usando la \textbf{negazione della doppia implicazione}: basta dimostrare che le negazioni dei due membri
    sono equivalenti.
    \begin{itemize}
        \item $f$ \textbf{non} è iniettiva se e solo se esistono $x_1 \neq x_2$ in $a$ con $f(x_1) = f(x_2)$
              (Proposizione~\ref{prop:forme-iniettivita}, punto 2).
        \item Negando il secondo membro: esiste $y \in b$ tale che $\overleftarrow{f}(\{y\})$ non è vuoto e non ha un solo
              elemento, cioè $\overleftarrow{f}(\{y\})$ ha \textbf{almeno due} elementi.
    \end{itemize}
    Queste due affermazioni sono equivalenti: se $x_1 \neq x_2$ e $f(x_1) = f(x_2)$, detto $y = f(x_1)$, l'insieme
    $\overleftarrow{f}(\{y\})$ contiene almeno $x_1$ e $x_2$; viceversa, se $\overleftarrow{f}(\{y\})$ contiene due elementi distinti
    $x_1$, $x_2$, allora $f(x_1) = y = f(x_2)$.
\end{dimostrazione}
\end{concetto}

\begin{concetto}
\begin{Teorema}{Caratterizzazione delle funzioni suriettive}{caratt-suriettive}
    \begin{prerequisiti}
        \dip{Definizione}{def:map-in-su-bi}
        \dip{Definizione}{def:immagine-antiimmagine}
        \dip{Definizione}{def:suriettiva}
        \dip{Proposizione}{prop:negdoppiaimp}
    \end{prerequisiti}
    \[ \forall a, b\ \forall f\,\Big(f \in \mathrm{Su}(a, b) \leftrightarrow \forall y \in b\,\big(\overleftarrow{f}(\{y\}) \neq \varnothing\big)\Big). \]
\end{Teorema}
\begin{dimostrazione}
    Di nuovo con la negazione della doppia implicazione: $f$ non è suriettiva se e solo se esiste $y \in b$ che non è
    immagine di nessun elemento di $a$, il che accade se e solo se $\overleftarrow{f}(\{y\}) = \{x \in a \mid f(x) = y\} = \varnothing$.
\end{dimostrazione}
\end{concetto}

\begin{concetto}
\begin{Corollario}{Caratterizzazione delle funzioni biettive con l'antiimmagine}{caratt-biettive-anti}
    \begin{prerequisiti}
        \dip{Teorema}{teo:caratt-iniettive}
        \dip{Teorema}{teo:caratt-suriettive}
        \dip{Definizione}{def:biettiva}
    \end{prerequisiti}
    $f \colon a \to b$ è biettiva se e solo se, per ogni $y \in b$, $\overleftarrow{f}(\{y\})$ ha esattamente un elemento, cioè
    \[ \forall y \in b\ \exists!\, x \in a\,\big(f(x) = y\big). \]
\end{Corollario}
\begin{dimostrazione}
    Si combinano i due teoremi precedenti: «al più un elemento» e «non vuoto» insieme danno «esattamente un elemento».
\end{dimostrazione}
\end{concetto}

% ──────────────────────────────────────────────────────────────
\section{Sezioni, retrazioni e inversa}\label{sec:sezioni}

\dfn[sezione-retrazione]{Sezione, retrazione, inversa}{
    Siano $f \colon a \to b$ e $g \colon b \to a$.
    \begin{itemize}
        \item Se $g \circ f = \mathrm{Id}_a$, allora $g$ si dice \textbf{retrazione} di $f$.
        \item Se $f \circ g = \mathrm{Id}_b$, allora $g$ si dice \textbf{sezione} di $f$.
        \item Se $g$ è sia sezione che retrazione, allora $g$ è l'\textbf{inversa} di $f$, e $f$ si dice \textbf{invertibile}.
    \end{itemize}
}

\begin{concetto}
\begin{Teorema}{Funzione iniettiva e retrazione}{iniettiva-retrazione}
    \begin{prerequisiti}
        \dip{Definizione}{def:sezione-retrazione}
        \dip{Definizione}{def:map-in-su-bi}
        \dip{Teorema}{teo:composte-caratteristiche}
        \dip{Definizione}{def:funzioni-notevoli}
        \dip{Definizione}{def:funzione}
    \end{prerequisiti}
    $f \colon a \to b$ è iniettiva $\leftrightarrow$ $a = \varnothing \ \vee\ \exists g \in \mathrm{Map}(b, a)\,(g \text{ è una retrazione di } f)$.
\end{Teorema}
\begin{dimostrazione}
    $(\rightarrow)$ Se $a = \varnothing$ non c'è niente da dimostrare. Altrimenti prendo $\bar{x} \in a \neq \varnothing$.
    Per ogni $y \in \mathrm{Im} f$ esiste un unico $x \in a$ con $f(x) = y$ (perché $f$ è iniettiva). Definisco
    \[ g \colon y \in b \mapsto \begin{cases} x & \text{se } y \in \mathrm{Im} f \wedge f(x) = y \\ \bar{x} & \text{se } y \in b \setminus \mathrm{Im} f \end{cases} \in a. \]
    Allora per ogni $x \in a$ si ha $g \circ f(x) = g(f(x)) = x$, cioè $g \circ f = \mathrm{Id}_a$: $g$ è una retrazione di $f$.

    $(\leftarrow)$ Se $a = \varnothing$, allora $f = (\varnothing, b, \varnothing)$ è iniettiva (la condizione di iniettività è
    un'implicazione con antecedente sempre falso). Se $a \neq \varnothing$, esiste $g \colon b \to a$ con $g \circ f = \mathrm{Id}_a$;
    l'identità è iniettiva, quindi $f$ è iniettiva per il Teorema~\ref{teo:composte-caratteristiche}, punto 1.
\end{dimostrazione}
\end{concetto}

\begin{concetto}
\begin{Teorema}{Funzione suriettiva e sezione}{suriettiva-sezione}
    \begin{prerequisiti}
        \dip{Definizione}{def:sezione-retrazione}
        \dip{Teorema}{teo:composte-caratteristiche}
    \end{prerequisiti}
    $f \colon a \to b$ è suriettiva $\leftrightarrow$ $f$ ha sezioni.
\end{Teorema}
\begin{dimostrazione}[Dimostrazione parziale]
    $(\leftarrow)$ Se $g$ è una sezione, $f \circ g = \mathrm{Id}_b$ è suriettiva, quindi $f$ è suriettiva per il
    Teorema~\ref{teo:composte-caratteristiche}, punto 2.
    $(\rightarrow)$ Di questa implicazione non si fornisce dimostrazione, perché richiederebbe l'introduzione di un ulteriore
    assioma: l'\textbf{Assioma della Scelta} (per ogni $y \in b$ bisogna «scegliere» un elemento di $\overleftarrow{f}(\{y\})$).
\end{dimostrazione}
\end{concetto}

\begin{concetto}
\begin{Teorema}{Sezioni e retrazioni}{unicita-inversa}
    \begin{prerequisiti}
        \dip{Definizione}{def:sezione-retrazione}
        \dip{Proposizione}{prop:assoc-prodotto}
        \dip{Definizione}{def:funzioni-notevoli}
    \end{prerequisiti}
    Sia $f \colon a \to b$. Sia $s$ una sezione di $f$ ed $r$ una retrazione di $f$. Allora $s = r$.
    In particolare, l'inversa, se esiste, è \textbf{unica}.
\end{Teorema}
\begin{dimostrazione}
    Usando l'associatività del prodotto (composizione):
    \[ r = r \circ \mathrm{Id}_b = r \circ (f \circ s) = (r \circ f) \circ s = \mathrm{Id}_a \circ s = s. \]
    Se $g$ e $g'$ sono due inverse, $g$ è una sezione e $g'$ una retrazione, quindi $g = g'$.
\end{dimostrazione}
\end{concetto}

Data la sua unicità, l'inversa di $f$, se esiste, la denoto con $f^{-1}$.

\begin{concetto}
\begin{Teorema}{Funzione biunivoca e inversa}{biettiva-inversa}
    \begin{prerequisiti}
        \dip{Teorema}{teo:iniettiva-retrazione}
        \dip{Teorema}{teo:suriettiva-sezione}
        \dip{Teorema}{teo:unicita-inversa}
        \dip{Definizione}{def:biettiva}
    \end{prerequisiti}
    $f \colon a \to b$ è biettiva $\leftrightarrow$ $f$ ha un'inversa.
\end{Teorema}
\begin{dimostrazione}
    È conseguenza dei teoremi precedenti.
    $(\leftarrow)$ L'inversa è una retrazione, quindi $f$ è iniettiva (Teorema~\ref{teo:iniettiva-retrazione}); è anche una
    sezione, quindi $f$ è suriettiva (Teorema~\ref{teo:suriettiva-sezione}).
    $(\rightarrow)$ $f$ suriettiva ha una sezione $s$. Se $a \neq \varnothing$, $f$ iniettiva ha una retrazione $r$;
    per il Teorema~\ref{teo:unicita-inversa} $s = r$, che è quindi l'inversa. Se $a = \varnothing$, la suriettività
    obbliga $b = \varnothing$ e l'inversa è la funzione vuota.
\end{dimostrazione}
\end{concetto}

\begin{concetto}
\begin{Teorema}{Funzione biettiva}{caratt-biettive}
    \begin{prerequisiti}
        \dip{Teorema}{teo:biettiva-inversa}
        \dip{Teorema}{teo:unicita-inversa}
        \dip{Corollario}{cor:caratt-biettive-anti}
        \dip{Teorema}{teo:suriettiva-sezione}
        \dip{Proposizione}{prop:forme-iniettivita}
        \dip{Definizione}{def:sezione-retrazione}
        \dip{Definizione}{def:immagine-antiimmagine}
    \end{prerequisiti}
    Sia $f \colon a \to b$. Le seguenti condizioni sono equivalenti:
    \begin{enumerate}
        \item $f$ è biettiva;
        \item $f$ ha sezioni e retrazioni;
        \item $f$ ha inversa;
        \item $f$ ha una sola sezione e una sola retrazione;
        \item $\forall y \in b\ \exists!\, x \in a\,\big(f(x) = y\big)$;
        \item $f$ ha una ed una sola sezione.
    \end{enumerate}
\end{Teorema}
\begin{dimostrazione}
    \textbf{(1) $\leftrightarrow$ (3)}: è il Teorema~\ref{teo:biettiva-inversa}.

    \textbf{(1) $\leftrightarrow$ (5)}: è il Corollario~\ref{cor:caratt-biettive-anti}.

    \textbf{(3) $\rightarrow$ (2)}: l'inversa è sia sezione sia retrazione.
    \textbf{(2) $\rightarrow$ (3)}: se $s$ è una sezione e $r$ una retrazione, per il Teorema~\ref{teo:unicita-inversa}
    $s = r$, che è quindi sia sezione sia retrazione, cioè l'inversa.

    \textbf{(3) $\rightarrow$ (4)}: sia $f^{-1}$ l'inversa. Ogni sezione $s$ è uguale alla retrazione $f^{-1}$, e ogni retrazione $r$
    è uguale alla sezione $f^{-1}$ (Teorema~\ref{teo:unicita-inversa}): quindi c'è una sola sezione e una sola retrazione.
    \textbf{(4) $\rightarrow$ (2)}: ovvio.

    \textbf{(3) $\rightarrow$ (6)}: come sopra, ogni sezione è uguale a $f^{-1}$.

    \textbf{(6) $\rightarrow$ (1)}: $f$ ha una sezione, quindi è suriettiva (Teorema~\ref{teo:suriettiva-sezione}).
    Per assurdo, $f$ non sia iniettiva: siano $x_1 \neq x_2$ elementi di $a$ con $f(x_1) = f(x_2) =: z$, e sia $s$ l'unica sezione di $f$.
    Dalla definizione di sezione $f(s(z)) = z$, cioè $s(z) \in \overleftarrow{f}(\{z\})$, che contiene almeno due elementi distinti.
    Posso quindi supporre, senza ledere la generalità, che $s(z) \neq x_1$ (altrimenti scambio $x_1$ e $x_2$). Definisco
    \[ g \colon y \in b \mapsto \begin{cases} s(y) & \text{se } y \neq z \\ x_1 & \text{se } y = z \end{cases} \in a. \]
    Per ogni $y \neq z$ si ha $f(g(y)) = f(s(y)) = y$, e $f(g(z)) = f(x_1) = z$: quindi $g$ è una sezione di $f$, diversa da $s$
    perché $g(z) = x_1 \neq s(z)$. Assurdo, perché la sezione era unica. Quindi $f$ è iniettiva, dunque biettiva.
\end{dimostrazione}
\end{concetto}
